\chapter{Regression tests}\label{ch:tests}
\module{NoCompromise.Tests.Constants,
        NoCompromise.Tests.Curvature,
        NoCompromise.Tests.Iteration,
        NoCompromise.Tests.LevelFunctionals,
        NoCompromise.Tests.Signs}

The tests below exist to catch sign, orientation and normalisation errors ---
the failure mode most likely to produce a Lean development that compiles and
proves something false.  Each should be an actual theorem in the Lean
development, not a comment.  They are grouped by which convention they probe.

\section{Conventions \ref{conv:outward} and \ref{conv:curvature}: normals and curvature}

\begin{test}[The model exterior potential]\label{test:model-potential}
For $u=\mathsf C/|x|$ on $\{|x|>r\}$, with the normal toward infinity,
\[
  H=\frac2r,
  \qquad
  w=\frac{\mathsf C}{r^2},
  \qquad
  \langle\nabla w,\nabla u\rangle=Hw^2 .
\]
A sign error in Convention~\ref{conv:curvature} flips the third identity.
\end{test}

\begin{test}[Stationarity of the round ball]\label{test:ball-stationary}
For $\Omega=B_R$, \eqref{eq:EL-pointwise} holds with
$H=2/R$ and $v_{B_R}=2\pi(R^2-r^2/3)$ from
\eqref{eq:ball-potential}; solving gives an explicit $\lambda$, which must
agree with \eqref{eq:scaling-identity}, i.e.\
$3V\lambda=2P_B+5D(B_V)$.  \emph{Checking both routes agree is the test.}
\end{test}

\begin{test}[The inner normal in Green's identity]\label{test:green-normal}
In \eqref{eq:green-second} the outward normal of the exterior on its inner
boundary is $-\nu_K$, so
\[
  \int_{\partial K}v_\Omega\,w\,d\Hs^2
  =-\int_{\partial K}\partial_{\nu_K}v_\Omega\,d\Hs^2 .
\]
Dropping the minus sign leaves \eqref{eq:green-identity} looking plausible but
gives $-V$.
\end{test}

\begin{test}[The tubular flux]\label{test:tubular-flux}
The tubular-cutoff computation of Lemma~\ref{lem:interior-flux} must return
\[
  \int_{\partial K}\partial_{\nu_K}v_\Omega\,d\Hs^2=-4\pi V ,
\]
\emph{not} $+4\pi V$ and not $-V$.  Test it on $\Omega=K=B_R$, where
$v_{B_R}(x)=V/|x|$ outside and $\partial_{\nu_K}v=-V/R^2$, giving
$-V/R^2\cdot4\pi R^2=-4\pi V$.
\end{test}

\section{Chapter \ref{ch:appK}: the level functionals}

\begin{test}[The model level functionals]\label{test:K-model}
For $u=\mathsf C/r$ the Euclidean level functions of
Chapter~\ref{ch:appK} satisfy
\[
  F(t)=4\pi t^2,
  \qquad
  p(t)=8\pi t,
  \qquad
  h(t)=0 .
\]
The third is the equality case of Lemma~\ref{lem:K-h-monotone} and is the
sharpest available check on \eqref{eq:K-h}.
\end{test}

\begin{test}[The Bochner density]\label{test:K-bochner}
For $u=\mathsf C/r$, $\Delta|\nabla u|=2\mathsf C/r^4$, matching
\eqref{eq:K-density} since a sphere of radius $r$ has
$|A|^2=2/r^2$ and $w=\mathsf C/r^2$, hence
$w|A|^2=2\mathsf C/r^4$ and $\nabla_\Sigma w=0$.
\end{test}

\section{Chapter \ref{ch:eps-reg}: the iteration}

\begin{test}[Normalisation of the recurrence]\label{test:recurrence}
The recurrence
\[
  E_{j+1}\le(\theta/2)E_j+C\theta^{j+1}\omega r
\]
must normalise, after division by $\theta^{j+1}$, to a geometric sum with
ratio $1/2$ (Lemma~\ref{lem:geometric-recurrence}).  If the ratio comes out
as $\theta/2$ or as $1$, the choice \eqref{eq:theta-choice} has been applied in
the wrong place.
\end{test}

\begin{test}[Excess of an exact halfspace and of a tilted plane]
\label{test:excess-values}
$\Exc(H,0,r,\nu)=0$ for the halfspace $H$ with outward normal $\nu$, and for a
halfspace with outward normal $\nu'$ at angle $\vartheta$ to $\nu$,
$\Exc(H',0,r,\nu)$ is a positive absolute multiple of $\vartheta^2$ for small
$\vartheta$.  This pins the normalisation $r^{-2}$ in
Definition~\ref{def:excess}.
\end{test}

\section{Chapter \ref{ch:flows}: the total curvature bound}

\begin{test}[Sphere and torus]\label{test:total-curvature}
For $\Sigma=\partial B_R$ with the outward normal, $\kappa\equiv1/R^2$, the
Gauss map is a diffeomorphism, $\iota_y(n)=1$ for every $y$, and
$\int_\Sigma\kappa\,d\Hs^2=4\pi$ --- equality in
\eqref{eq:total-curvature-bound}.  For a standard torus of revolution,
$\int_\Sigma\kappa\,d\Hs^2=0$, and a height function in a generic direction has
$c_0=c_2=1$, $c_1=2$, so
$c_0-c_1+c_2=0$: this exercises
Proposition~\ref{prop:index-sum-antipodal} and
Proposition~\ref{prop:morse-count} in a case where the inequality is slack and
$\kappa$ changes sign.
\end{test}

\begin{test}[The antipodal averaging is not vacuous]\label{test:antipodal}
On the round sphere, $\iota_y(n)=\iota_{-y}(n)=1$, so
$\iota_y+\iota_{-y}=2$ and Proposition~\ref{prop:morse-count} is tight.  A
version of \eqref{eq:index-sum-antipodal} that dropped the antipodal term ---
bounding $\iota_y(n)$ alone by $2$ rather than $\iota_y+\iota_{-y}$ --- would
still pass on the sphere but would give only
$\int_\Sigma\kappa\le8\pi$, which is too weak for
Chapter~\ref{ch:appK}.  Test the factor $\tfrac12$ explicitly.
\end{test}

\begin{test}[Curvature sign convention]\label{test:kappa-sign}
For $\Sigma=\partial B_R$ with the \emph{outward} normal, the shape operator
of Definition~\ref{def:frame-free} is $A=\tfrac1RI$, so $H=2/R>0$ and
$\kappa=1/R^2>0$, matching Convention~\ref{conv:curvature}.  With the inward
normal both signs of $H$ flip while $\kappa$ does not --- which is why
\eqref{eq:sign-kappa-index} is orientation-independent and
\eqref{eq:hess-height} is not.
\end{test}

\section{Chapter \ref{ch:cones}: the product competitor}

Regression test~\ref{test:product-cost} (horizontal cost exactly
$2|A\triangle L|$) belongs here.

\section{Chapter \ref{ch:algebra}: the constants}

\begin{test}[Headline constants]\label{test:constants}
The following must be provable numerically (or as exact algebraic identities):
\[
  V_*=\frac{5\cdot2^{2/3}}{2^{1/3}+1}\approx3.51207,
  \qquad
  \frac52<V_*<4,
\]
the optimal ball volume is $5/2$, and
$e_*=3(9\pi/5)^{1/3}$.  The ratio of the two ways of writing $V_*$ in
Definition~\ref{def:threshold} and \eqref{eq:Vstar-form} must be exactly $1$.
\end{test}

